% Additive section. Does not modify the 211-page edition.
% Provenance and checks: internal documents in this same directory.
\section{Nonadic refinement, positive flags, and compatibility of moments}
\label{sec:dim-cofinal}

The length of a finite representation determines how many minors can be
formed from its columns. Refinement preserves inherited columns and
adds others with known provenance. Thus the dimension of a matrix
realization can grow without replacing the register that fixes its
separations. The following construction combines this growth with two
precise notions of compatibility: restriction of evaluations to the
previous marks and averaging of observables over descendant cells.
The former preserves polynomial flags; the latter preserves a
measure and separates the information increment into orthogonal details.

The starting datum is the dodecaphasic output already produced by
APP--TRIT--TPK. The two APP sheets preserve their complete divisions;
TRIT orients the regime, and the operator
\(U_t=\operatorname{Upd}_t\circ\operatorname{Tra}_t\circ
\operatorname{Sel}_t\) transports the enriched state. The prolongation
\(w_6\to w_{12}\to w_{18}\to w_{24}\to w_{30}\to R_{36}\to G_9\)
preserves phase return and memory advance. The joint discrete structure
of the continuum correlatively preserves the representations of that state,
with their sheets, orientation, memory, and boundaries. The regional
representations and the signed register come from this same state;
inversion of the latter gives
\[
 K=(234,543,140,729,659,824,621,58,914,794,146,601),
 \qquad Q=\sum_{j=1}^{12}K_j=6263.
\]
This output is used with its marks and orientation; the moment matrices
are subsequent realizations. The refinement index introduced below
counts subdivisions of the geometric reader. Identifying it with a
physical duration or a number of TPK microtransitions requires preserving
the clock of that other realization.

\subsection{Marked partitions and cofinal growth}

Write \(d_j=K_j/Q\), \(a_0=0\), and
\(a_j=\sum_{h=1}^j d_h\). For \(r\geq0\), define
\[
 L_r=12\,9^r,\qquad N_r=L_r+1,\qquad
 t_{r,(j-1)9^r+c}=a_{j-1}+\frac{c}{9^r}d_j,
 \quad 0\leq c\leq9^r.
\]
Shared endpoints are counted only once. The \(L_r\) cells are indexed
by \((j,c)\), with \(1\leq j\leq12\) and \(0\leq c<9^r\),
and have length \(\ell_{r,j,c}=d_j/9^r\). The expression of \(c\)
in \(r\) nonadic digits, including leading zeros, is also preserved.
A successor has index \(9c+h\), with \(0\leq h\leq8\); Euclidean
division recovers its predecessor and last digit. Shared endpoints retain
their boundary mark; the half-open interval convention only organizes
the cells and removes no mark.

\begin{proposition}[Cofinality and recovery of the register]
\label{prop:dim-grid}
The marks satisfy
\[
 0=t_{r,0}<t_{r,1}<\cdots<t_{r,L_r}=1,
 \qquad t_{r+1,9i}=t_{r,i}.
\]
The maximum length is \(h_r=(\max_jK_j)/(Q9^r)\), which tends
to zero. For every pair of integers \(k\geq1\), \(m\geq1\),
there is an \(r_0\) such that \(k+m\leq N_r\) for \(r\geq r_0\).
The register is recovered at any level through
\[
 K_j=Q\sum_{c=0}^{9^r-1}\ell_{r,j,c}.
\]
\end{proposition}
\begin{proof}
Each consecutive difference is \(d_j/9^r>0\). Replacing \(c\)
by \(9c\) in the next-level formula gives exactly the previous mark.
The nine successor lengths sum to the predecessor length; iteration proves
recovery. Finally, \(N_r=12\,9^r+1\) grows without bound.
An explicit choice is any \(r_0\geq0\) with
\(9^{r_0}\geq(k+m-1)/12\). The entire argument preserves \(Q\),
the block index, and the descent digits.
\end{proof}

Normalization determines the ratios \(K_j/Q\). Recovery of the
integers also uses the charge \(Q\) preserved by the register.
This distinguishes the normalized geometry from the scale information
accompanying its representation.

\subsection{Evaluation flags and realizations in every finite dimension}

For \(1\leq q\leq N_r\), let
\[
 V_r^{(q)}=(t_{r,i}^{\,a})_{0\leq a<q,\,0\leq i<N_r},
 \qquad F_r^q=\operatorname{row}V_r^{(q)}.
\]
The zeroth power is one at the zero node as well.

\begin{theorem}[Positive flag and arbitrary-rank composition]
\label{thm:dim-all-ranks}
The spaces \(F_r^q\) form a flag
\(F_r^1\subset F_r^2\subset\cdots\subset F_r^{N_r}\), with
\(\dim F_r^q=q\), and every matrix \(V_r^{(q)}\) has all its
ordered maximal minors strictly positive. For
\(k\geq1\), \(m\geq1\), \(k+m\leq N_r\), the matrices
\[
 C_r=V_r^{(k)},\quad Z_r=V_r^{(k+m)},\quad
 Y_r=C_rZ_r^{\mathsf T}
\]
satisfy
\[
 (Y_r)_{ab}=\sum_{i=0}^{L_r}t_{r,i}^{a+b},\qquad
 \operatorname{rank}Y_r=k,\qquad
 [Y_r]\in\mathcal A_{N_r,k,m}(Z_r).
\]
The construction reaches every finite pair \((k,m)\) through a
sufficiently deep level of the same marked partition.
\end{theorem}
\begin{proof}
For \(i_1<\cdots<i_q\), the corresponding minor is
\[
 \det(t_{r,i_b}^{a-1})_{a,b=1}^{q}
   =\prod_{a<b}(t_{r,i_b}-t_{r,i_a})>0.
\]
This proves rank, positivity, and the flag inclusions. The initial
\(k\times k\) block of \(Y_r\) is the Gram matrix of the evaluations
of \(1,t,\ldots,t^{k-1}\). For \(x\neq0\),
\[
 x^{\mathsf T}(Y_r)_{[0,k-1],[0,k-1]}x
 =\sum_i\left(\sum_{a=0}^{k-1}x_at_{r,i}^{a}\right)^2>0,
\]
because a nonzero polynomial of degree less than \(k\) cannot vanish
at the \(N_r\geq k+m\) distinct nodes. The block is invertible,
and \(Y_r\), which has \(k\) rows, has rank \(k\).
The stated membership is the application of the positive external
matrix \(Z_r\) to the positive point represented by \(C_r\).
\end{proof}

We use the mathematical convention
\[
 \mathcal A_{N,k,m}(Z)=
 \{[CZ^{\mathsf T}]:[C]\in\operatorname{Gr}_{\geq0}(k,N)\}.
\]
The definition for general \(m\) can be found in S.~N.~Karp,
L.~K.~Williams, and Y.~X.~Zhang,
\href{https://people.math.harvard.edu/~williams/papers/amplituhedron-plus.pdf}
{\emph{Decompositions of amplituhedra}}, Definition~1.1.
Here the external datum already comes from the partition expressed by HMT.
The conclusion constructs a determined collection of points and flags;
the dimensions of the ambient Grassmannian and of an atlas of its
image are objects additional to the ranks proved.

Let \(E_r:\mathbb R^{N_{r+1}}\to\mathbb R^{N_r}\) be restriction
to the positions \(9i\). For every \(q\leq N_r\),
\[
 E_r\bigl((V_{r+1}^{(q)})^{\mathsf T}x\bigr)
       =(V_r^{(q)})^{\mathsf T}x.
\]
Thus \(E_r:F_{r+1}^q\to F_r^q\) is an isomorphism on these
evaluation spaces: its two realizations identify the same polynomial
of degree less than \(q\). Composing restrictions agrees with direct
restriction. Minors on inherited indices preserve their value exactly.

\begin{proposition}[Recovery from a calibrated flag]
If the column order, endpoint marks, and calibration
\(t_{r,0}=0\), \(t_{r,L_r}=1\) are preserved, the plane \(F_r^2\)
determines the nodes and separations. Together with \(Q\) and the
descent marks, it determines \(K\).
\end{proposition}
\begin{proof}
The plane contains the constant vector and the node vector. Any other
coordinate that, together with the constant vector, generates the same
plane is \(a+bt\), with \(b\neq0\). The two calibrations fix
\(a=0\), \(b=1\). Consecutive differences recover the lengths,
and Proposition~\ref{prop:dim-grid} recovers the register.
\end{proof}

\subsection{Rank-one convex bodies and growth of the domain}

For \(m\geq1\) and \(m+1\leq N_r\), define
\[
 P_r^{(m)}=\operatorname{conv}
 \{(t_{r,i},t_{r,i}^2,\ldots,t_{r,i}^m):0\leq i\leq L_r\}.
\]
Normalizing a positive row of \(C\) as barycentric coefficients proves
\(\mathcal A_{N_r,1,m}(Z_r)=P_r^{(m)}\). The Vandermonde
determinant provides \(m+1\) affinely independent points, so the
body has dimension \(m\).

\begin{proposition}[Simplicial covering in finite dimension and convex limit]
\label{prop:dim-polytopes}
Every \(P_r^{(m)}\) admits a triangulation using its vertices, with
degree one over the interior of every simplicial cell. The sum of
their canonical forms has canceled interior residues. Moreover,
\[
 P_r^{(m)}\subseteq P_{r+1}^{(m)},\qquad
 \overline{\bigcup_{r\geq r_0}P_r^{(m)}}=
 \operatorname{conv}\{(t,t^2,\ldots,t^m):0\leq t\leq1\}.
\]
\end{proposition}
\begin{proof}
Order the vertices by their marks. In each face, choose its least
vertex; recursively triangulate the facets not containing it and form
the corresponding cones. Induction on dimension proves covering and
disjoint interiors, and the same rule on shared faces guarantees
compatibility. In each oriented simplex \([v_0,\ldots,v_m]\), with
matrix \(A=(v_1-v_0,\ldots,v_m-v_0)\) of positive determinant and
barycentric coordinates \(\lambda_0,\ldots,\lambda_m\), the
barycentric map is affine and invertible. Its form is
\[
 \Omega_{[v_0,\ldots,v_m]}=
 \frac{dx_1\wedge\cdots\wedge dx_m}
 {\det A\prod_{j=0}^{m}\lambda_j}.
\]
The residual restriction to each face is its oriented form. Two
adjacent simplices induce opposite orientations on their common face,
so the interior residues cancel. This is the same triangulation and
residue operation applied in the previous rank-one realization, with
the number of vertices now increased by refinement.

Inclusion of inherited nodes proves inclusion of the bodies. The mesh
tends to zero, so the sampled curves are dense in the compact moment
curve. Every finite convex combination of curve points is approximated
by combinations using nodes from a common level. Conversely, all the
bodies lie in the compact convex hull of the curve. Both inclusions
prove the equality.
\end{proof}

Cancellation of residues acts on a triangulation or subdivision of a
fixed body. Passing from \(P_r^{(m)}\) to \(P_{r+1}^{(m)}\)
adds vertices and can enlarge the body. Their forms belong to their
respective domains. The preceding convex limit specifies convergence
of the bodies without imposing an identity between their forms.

\subsection{Calendar measure and convergence of nodal moments}

Let \(F_K:[0,1]\to[0,1]\) be the increasing piecewise affine
function taking \((j-1)/12\) to \(a_{j-1}\) and \(j/12\) to
\(a_j\). On the corresponding interval, its slope is \(12d_j>0\).
The transported calendar measure is
\[
 \mu_K=(F_K)_*(ds),\qquad
 \frac{d\mu_K}{dt}=\frac1{12d_j}\quad
 (a_{j-1}<t<a_j).
\]
Each initial interval has mass \(1/12\), and each level-\(r\)
cell has mass \(1/L_r\). This choice assigns equal weight to the
calendar positions before realizing their lengths. The length measure
\(dt\) is another reader; the two coincide when all initial
separations are equal.

Its moments are rational and determined by \(K\):
\begin{equation}
 M_p=\int_0^1t^p\,d\mu_K(t)
   =\frac1{12(p+1)}\sum_{j=1}^{12}
       \frac{a_j^{p+1}-a_{j-1}^{p+1}}{d_j},\qquad p\geq0.
 \label{eq:dim-moments}
\end{equation}
\begin{proposition}[Recovery through moments and quantiles]
\label{prop:dim-moment-recovery}
The full sequence \((M_p)_{p\geq0}\) determines \(\mu_K\).
If \(H_K(t)=\mu_K([0,t])\) is its distribution function, then
\[
 a_j=H_K^{-1}(j/12),\qquad
 K_j=Q\bigl(H_K^{-1}(j/12)-H_K^{-1}((j-1)/12)\bigr).
\]
Consequently, the complete moments, the charge \(Q\), and the calendar
order recover the initial register.
\end{proposition}
\begin{proof}
Two probability measures on \([0,1]\) with the same moments have the
same integrals for all polynomials. Uniform approximation of continuous
functions by polynomials and finite mass extend this equality to all
continuous functions, which determine the Borel measure. The density
of \(\mu_K\) is positive and constant on each initial interval.
Its distribution is therefore continuous and strictly increasing, and
satisfies \(H_K(a_j)=j/12\). Inverting this equality and taking the
difference of endpoints prove the formulas.
\end{proof}

The nodal probability measure corresponding to Theorem
\ref{thm:dim-all-ranks} is
\(\nu_r=N_r^{-1}\sum_{i=0}^{L_r}\delta_{t_{r,i}}\).

\begin{theorem}[Limit of the nodal representation]
\label{thm:dim-nodal-limit}
We have \(\nu_r\rightharpoonup\mu_K\). If \(f\) is continuous,
\[
 \left|\int f\,d\nu_r-\int f\,d\mu_K\right|
 \leq \omega_f(h_r)+\frac{\operatorname{osc}(f)}{N_r},
\]
where \(\omega_f\) is its modulus of continuity. For \(p\geq1\),
the error of the moment of order \(p\) is bounded by
\(p h_r+1/N_r\). For every fixed pair \((k,m)\),
\[
 \frac1{N_r}Y_r\longrightarrow (M_{a+b})_{
     0\leq a<k,\,0\leq b<k+m}=:Y_\infty^{(k,m)},
 \qquad\operatorname{rank}Y_\infty^{(k,m)}=k.
\]
\end{theorem}
\begin{proof}
The measure assigning mass \(1/L_r\) to each left cell endpoint
approximates its integral with error at most \(\omega_f(h_r)\).
Indeed, each cell has that mass and diameter no greater than \(h_r\).
The measure \(\nu_r\) is \(L_r/N_r\) times this measure plus
\(\delta_1/N_r\), adding at most
\(\operatorname{osc}(f)/N_r\). For \(f(t)=t^p\), the Lipschitz
constant is \(p\). Convergence of each entry follows from these
bounds. The initial block of the limit matrix is positive definite:
the integral of the square of a nonzero polynomial is positive,
because \(\mu_K\) has strictly positive density on every open
interval of \([0,1]\). Thus the limit preserves rank \(k\).
\end{proof}

Nodal moments change exactly through the addition of nodes.
If \(\mathcal J_r\) is the set of new positions and
\(z_q(t)=(1,t,\ldots,t^{q-1})^{\mathsf T}\), then
\[
 V_{r+1}^{(q)}(V_{r+1}^{(q)})^{\mathsf T}
 =V_r^{(q)}(V_r^{(q)})^{\mathsf T}
  +\sum_{i\in\mathcal J_r}z_q(t_{r+1,i})z_q(t_{r+1,i})^{\mathsf T}.
\]
Each summand is positive semidefinite. The same update holds for the
rectangular matrices \(Y_r\), using \(z_kz_{k+m}^{\mathsf T}\).
The equality expresses the change in counting measure; normalization
and the preceding theorem describe its limit. The rank-\(k\) limit
object has coordinates in a finite Grassmannian, whereas the external
data \(Z_r\) belong to the different representations at each level.

\subsection{Cell averages, compatible operators, and orthogonal detail}

For a cell \(I_{r,i}=[b_{r,i},b_{r,i}+\ell_{r,i}]\), set
\[
 a_{r,i}^{(p)}=
 \frac{(b_{r,i}+\ell_{r,i})^{p+1}-b_{r,i}^{p+1}}
 {(p+1)\ell_{r,i}},\qquad
 A_r^{(q)}=(a_{r,i}^{(p)})_{0\leq i<L_r,\,0\leq p<q}.
\]
These are the monomial means with respect to the conditional probability
of \(\mu_K\) on each cell. This probability is uniform on the cell,
although \(\mu_K\) is not the global length measure.

Let \(\mathcal H_r=\mathbb R^{L_r}\) with inner product
\(\langle x,y\rangle_r=L_r^{-1}\sum_i x_i y_i\). The map
\(J_r:\mathcal H_r\to\mathcal H_{r+1}\) repeats a predecessor's value
on its nine successors. Its adjoint \(R_r=J_r^*\) averages those
nine values. We have \(R_rJ_r=I\).

\begin{theorem}[Exact compatibility and moment balance]
\label{thm:dim-cell-balance}
For every \(p\geq0\),
\[
 R_r a_{r+1}^{(p)}=a_r^{(p)},\qquad
 \frac1{L_r}\sum_i a_{r,i}^{(p)}=M_p.
\]
Define
\[
 D_r^{(q)}=A_{r+1}^{(q)}-J_rA_r^{(q)},\quad
 G_r^{(q)}=\frac1{L_r}(A_r^{(q)})^{\mathsf T}A_r^{(q)}.
\]
Then
\[
 R_rD_r^{(q)}=0,\qquad
 G_{r+1}^{(q)}=G_r^{(q)}+
       \frac1{L_{r+1}}(D_r^{(q)})^{\mathsf T}D_r^{(q)},
\]
and
\[
 G_r^{(q)}\longrightarrow H_K^{(q)}=(M_{a+b})_{0\leq a,b<q}.
\]
The details and initial level reconstruct all levels. For
\(q\leq L_r\), the matrix \(A_r^{(q)}\) has rank \(q\), and
all ordered minors of order \(q\) are positive.
\end{theorem}
\begin{proof}
The nine successors have equal masses, and their intervals partition the
predecessor interval. Additivity of the integral gives the first identity;
summing over all cells gives the second. The decomposition
\(A_{r+1}=J_rA_r+D_r\) is orthogonal in each column because
\(R_rD_r=0\). Expanding its Gram matrix cancels the cross terms,
and \(J_r\) is an isometry. This proves the matrix balance.

The cellwise constant observable taking the value \(a_{r,i}^{(p)}\)
uniformly approximates \(t^p\), with error no greater than \(p h_r\)
for \(p\geq1\). It therefore converges in \(L^2(\mu_K)\), and
the inner products converge to \(M_{a+b}\). Reconstruction is
obtained by iterating \(A_{r+1}=J_rA_r+D_r\).

To prove the minors, choose \(q\) ordered cells. By multilinearity
and Fubini, the determinant of their averages is the mean over the
product of these cells of \(\prod_{a<b}(x_b-x_a)\). This integrand
is positive in the interior of the product because the cells are
disjoint and ordered. Their boundaries have measure zero. The mean
is strictly positive.
\end{proof}

The identity extends to the degree pair \(k\) and \(k+m\):
\[
 \overline Y_r=\frac1{L_r}(A_r^{(k)})^{\mathsf T}A_r^{(k+m)},
 \qquad
 \overline Y_{r+1}-\overline Y_r
 =\frac1{L_{r+1}}(D_r^{(k)})^{\mathsf T}D_r^{(k+m)}.
\]
For \(k+m\leq L_r\), the averaging matrices define a second
positive amplituhedral realization, with \(L_r\) columns, and
\(\overline Y_r\) has rank \(k\). These are cell averages,
whereas \(Y_r\) uses evaluations at \(N_r=L_r+1\) nodes. Both
realizations converge to \(Y_\infty^{(k,m)}\); their agreement
in the limit follows from the preceding proofs.

The operator version is equally explicit. For a continuous function
\(f\), let \(\Psi_r(f)\) be the diagonal operator of its conditional
means on \(\mathcal H_r\). It is linear, positive, and unital, and
\[
 R_r\Psi_{r+1}(f)J_r=\Psi_r(f).
\]
This equality is proved by applying both sides to each basis vector.
For \(X(t)=t\), the first two operators determine, cell by cell,
\[
 \Psi_r(X)_i=b_{r,i}+\frac{\ell_{r,i}}2,\qquad
 \Psi_r(X^2)_i-\Psi_r(X)_i^2=\frac{\ell_{r,i}^2}{12}.
\]
If \(v_{r,i}\) denotes this variance, the formula
\(\ell_{r,i}=\sqrt{12v_{r,i}}\) recovers the length; the first
identity recovers the initial endpoint. With orientation, marks, and
\(Q\), one again recovers \(K\). Averaging is a positive compression:
the displayed variance quantifies why \(\Psi_r(X^2)\) and
\(\Psi_r(X)^2\) are different operations.

\subsection{Transport of the measure and scope of dimensional growth}

Two positive lists of separations \(d\) and \(d'\), with the same
marks, determine an increasing piecewise affine homeomorphism
\(F=F_{K'}\circ F_K^{-1}\). It takes each cell to its corresponding
cell and satisfies \(F_*\mu_K=\mu_{K'}\). Hence
\[
 \mathcal U_F:L^2(\mu_K)\longrightarrow L^2(\mu_{K'}),
 \qquad\mathcal U_F f=f\circ F^{-1}
\]
is a surjective isometry. If \(P_r\) and \(P'_r\) are the
projections onto cellwise constant functions, then
\(P'_r\mathcal U_F=\mathcal U_FP_r\): conditional means are
transported because the masses of corresponding cells agree.
This applies in particular to the positive separation flow directed
by \(K\), preserving each block's direction in its successors. The
moments of the spatial coordinate may change; the isometry transports
the observable together with the measure and its projections.

The growth proved is cofinal in every finite external dimension.
It preserves a system of positive flags, an explicit collection of
amplituhedral points, complete convex bodies in rank one, and a measure
whose averaging operators are exactly compatible. Memory of the details
allows inversion of compression between levels. The rank-one covering
and residue proofs preserve their domain; for higher ranks, the present
result fixes matrices, positivity, recovery, compatibility, and limit
with explicit types. This separation locates the scope of each
construction within the same development without confusing dimension
growth, parametrization of a point, and classification of all cells
of an image.
